% =====================================================================
\section{要点}

この章では、フォルダの地図 → データの形 → 1つの機能の流れ → 細部、の順に読んでいく。

\subsection{フォルダの地図}

\begin{kkbadgepoint}[要点]{最初に読むのは赤の場所だけでいい}
  \texttt{mobile/} には約50個のファイルがあるが、このアプリ独自の中心は
  \texttt{app/(tabs)/}・\texttt{context/}・\texttt{types/} の3か所。
  \texttt{components/} の一部や \texttt{hooks/} の一部は、Expo が最初に用意した雛形がそのまま残っているだけで、
  4つは今どこからも使われていない。自動生成の \texttt{node\_modules/} や \texttt{dist/} は読む必要がない。
\end{kkbadgepoint}

\textgt{\texttt{mobile/}（アプリ本体）の中身}

\begin{longtblr}[entry=none,label=none]{
  colspec={Q[l,wd=4.3cm] X[l] Q[c,wd=1.8cm] Q[c,wd=1.6cm]},
  width=\linewidth,
  rowhead=1,
  row{1}={bg=kkdark, fg=white, font=\small\gtfamily},
  hlines={black!15},
  cells={font=\small},
}
場所 & 役割（一言） & 種類 & 重要度 \\
\SetRow{bg=impAbg} \texttt{app/(tabs)/\_layout.tsx} & 下のタブ5つと、共有データの箱（Provider）の配置 & 独自 & \lvA \\
\SetRow{bg=impAbg} \texttt{app/(tabs)/index.tsx} & ホーム（キャラ・累計・今日の時間） & 独自 & \lvA \\
\SetRow{bg=impAbg} \texttt{app/(tabs)/timer.tsx} & タイマー（開始・停止と前後チェック） & 独自 & \lvA \\
\SetRow{bg=impBbg} \texttt{app/(tabs)/history.tsx} & カレンダー（日ごとの学習時間） & 独自 & \lvB \\
\SetRow{bg=impBbg} \texttt{app/(tabs)/fatigue.tsx} & 疲労チェック（予定の入力とバッテリー） & 独自 & \lvB \\
\SetRow{bg=impBbg} \texttt{app/(tabs)/settings.tsx} & 設定（テーマ・換算・グラフ・書き出し・削除） & 独自 & \lvB \\
\SetRow{bg=impBbg} \texttt{app/\_layout.tsx} & 一番外側の枠。ログインしていなければログイン画面へ送る & 独自 & \lvB \\
\SetRow{bg=impBbg} \texttt{app/login.tsx} & ログイン・新規登録・ゲストで使う & 独自 & \lvB \\
\SetRow{bg=impCbg} \texttt{app/modal.tsx} & 雛形の見本画面。どこからも開かれていない & 雛形（未使用） & \lvC \\
\SetRow{bg=impAbg} \texttt{context/} & アプリ全体で共有するデータと、その保存・読込（5つ） & 独自 & \lvA \\
\SetRow{bg=impAbg} \texttt{types/} & データの形（型）の定義だけ（2つ） & 独自 & \lvA \\
\SetRow{bg=impBbg} \texttt{utils/} & 計算・変換の関数。画面を持たない（7つ） & 独自 & \lvB \\
\SetRow{bg=impBbg} \texttt{hooks/use-fatigue.ts} & 疲労の状態を計算して画面に渡す & 独自 & \lvB \\
\SetRow{bg=impCbg} \texttt{hooks/} の残り3つ & ダーク／ライトの判定と色選び & 雛形 & \lvC \\
\SetRow{bg=impBbg} \texttt{components/} の独自8つ & 疲労バッテリー・ぼかし・反応時間テスト・点数選び・グラフ2つ・前後チェック2つ & 独自 & \lvB \\
\SetRow{bg=impCbg} \texttt{components/} の雛形4つ（使用中） & \texttt{themed-text}・\texttt{themed-view}・\texttt{haptic-tab}・\texttt{ui/icon-symbol} & 雛形 & \lvC \\
\SetRow{bg=impDbg} \texttt{components/} の雛形4つ（未使用） & \texttt{hello-wave}・\texttt{parallax-scroll-view}・\texttt{external-link}・\texttt{ui/collapsible} & 雛形（未使用） & \lvD \\
\SetRow{bg=impBbg} \texttt{lib/supabase.ts} & Supabase（ログイン用サーバー）への接続 & 独自 & \lvB \\
\SetRow{bg=impCbg} \texttt{constants/theme.ts} & 色とフォント。\texttt{Colors} は雛形、\texttt{AppColors}（緑）は独自 & 混在 & \lvC \\
\SetRow{bg=impBbg} \texttt{package.json} & 使うライブラリとコマンドの一覧 & 設定 & \lvB \\
\SetRow{bg=impCbg} \texttt{app.json} & アプリ名・アイコン・対応 OS など Expo の設定 & 設定 & \lvC \\
\SetRow{bg=impCbg} \texttt{tsconfig.json} & TypeScript の設定（\texttt{strict}、\texttt{@/} の別名） & 設定 & \lvC \\
\SetRow{bg=impCbg} \texttt{eslint.config.js} & 書き方のチェックの設定 & 雛形 & \lvC \\
\SetRow{bg=impCbg} \texttt{assets/} & アイコンや起動画面の画像 & 素材 & \lvC \\
\SetRow{bg=impCbg} \texttt{README.md} & Expo の雛形の説明書のまま & 雛形 & \lvC \\
\SetRow{bg=impDbg} \texttt{node\_modules/}・\texttt{package-lock.json} & インストールしたライブラリの本体と、その版の記録 & 自動生成 & \lvD \\
\SetRow{bg=impDbg} \texttt{.expo/}・\texttt{dist/}・\texttt{expo-env.d.ts} & 開発中の一時ファイル・Web 版の完成品・型の補助 & 自動生成 & \lvD \\
\SetRow{bg=impDbg} \texttt{.env.local} & Supabase の接続先などの秘密の値。開かない & 秘密 & \lvD \\
\end{longtblr}

\texttt{@/} は \texttt{tsconfig.json} で決めた「\texttt{mobile/} を指す別名」。
\texttt{import \{ useTime \} from '@/context/TimeContext'} は「\texttt{mobile/context/TimeContext.tsx} から読み込む」という意味になる。

\textgt{リポジトリ直下（アプリの外）}

\begin{center}
\begin{tblr}{
  colspec={Q[l,wd=4.3cm] X[l] Q[c,wd=1.6cm]},
  width=\linewidth,
  row{1}={bg=kkdark, fg=white, font=\small\gtfamily},
  hlines={black!15},
  cells={font=\small},
}
場所 & 役割（一言） & 重要度 \\
\SetRow{bg=impBbg} \texttt{log/} & 開発ログ V1〜V3。「なぜ」の手がかりがある & \lvB \\
\SetRow{bg=impCbg} \texttt{CLAUDE.md}・\texttt{AGENTS.md} & AI への指示（フォルダ構成と安全ルール） & \lvC \\
\SetRow{bg=impCbg} \texttt{netlify.toml} & Web 版を公開する設定 & \lvC \\
\SetRow{bg=impCbg} \texttt{.gitignore} & Git に入れないファイルの一覧 & \lvC \\
\end{tblr}
\end{center}

\subsection{データの形（types/）}

\begin{kkbadgepoint}[要点]{コードを読む前に「何を覚えているアプリか」を知る}
  \texttt{types/} には処理が1行もなく、「データがどんな項目でできているか」だけが書いてある。
  ここを先に見ると、画面や context のコードが「この項目を作っている」「この項目を足している」と読めるようになる。
  学習時間そのもの（累計と日ごと）は \texttt{types/} ではなく、\texttt{context/TimeContext.tsx} の中で形を決めている。
\end{kkbadgepoint}

\textgt{TypeScript の型の読み方}（このあとのコードに出てくるもの）
\begin{itemize}
  \item \texttt{interface 名前 \{ … \}}：項目の一覧。Python の \texttt{TypedDict} や \texttt{dataclass} に近い。
  \item \texttt{subject?: string}：\texttt{?} は「あってもなくてもよい」。Python の \texttt{Optional[str]}。
  \item \texttt{'pre' | 'post'}：\texttt{|} は「どれか1つ」。この場合は2つの文字列のどちらか。Python の \texttt{Literal['pre', 'post']}。
  \item \texttt{number | null}：数字か、何もない（\texttt{null}）か。
  \item \texttt{Record<string, number>}：キーが文字列、値が数字の辞書。Python の \texttt{dict[str, int]}。
  \item \texttt{ReactionTrial[]}：\texttt{[]} は配列。Python の \texttt{list[ReactionTrial]}。
\end{itemize}

\begin{center}
\begin{tikzpicture}[font=\footnotesize,
  ent/.style={draw=#1, fill=white, line width=0.9pt, rounded corners=2pt, inner sep=4pt},
  arr/.style={-{Stealth[length=2.2mm]}, line width=0.8pt, draw=black!60},
  lab/.style={font=\scriptsize, fill=white, inner sep=1.2pt}]
  \node[ent=impA] (ss) at (0,0) {\begin{tabular}{@{}l@{}}
    \textcolor{impA}{\textgt{StudySession}}（1回の学習）\\\hline
    \texttt{id}：この記録の番号\\
    \texttt{subject?}：科目\\
    \texttt{startedAt}／\texttt{endedAt?}：開始・終了の日時\\
    \texttt{durationMinutes}：学習した分数\\
    \texttt{reward}：面白さ・重要度・締切\\
    \texttt{subjectivePreId?} ほか：チェックの番号
  \end{tabular}};
  \node[ent=impB] (se) at (7.9,1.3) {\begin{tabular}{@{}l@{}}
    \textcolor{impB}{\textgt{SubjectiveEntry}}（自分で感じた疲れ）\\\hline
    \texttt{sessionId?}：どの学習のものか\\
    \texttt{timing}：\texttt{'pre'}（前）か \texttt{'post'}（後）\\
    \texttt{score}：疲労感 1〜5\\
    \texttt{note?}：メモ
  \end{tabular}};
  \node[ent=impB] (rr) at (7.9,-2.1) {\begin{tabular}{@{}l@{}}
    \textcolor{impB}{\textgt{ReactionTimeResult}}（反応時間テスト）\\\hline
    \texttt{sessionId?}：どの学習のものか\\
    \texttt{timing}：前か後\\
    \texttt{trials}：1回ずつの結果\\
    \texttt{avgReactionMs}：平均の反応時間\\
    \texttt{accuracy}：正確さ 0〜1
  \end{tabular}};
  \node[ent=impC] (ra) at (0,-4.3) {\begin{tabular}{@{}l@{}}
    \textcolor{impC}{\textgt{RewardAppraisal}}（やる気の材料）\\\hline
    \texttt{interest}：面白さ 1〜5\\
    \texttt{importance}：重要度 1〜5\\
    \texttt{hasDeadline}／\texttt{deadlineUrgency?}：締切
  \end{tabular}};
  \node[ent=impC] (rt) at (7.9,-5.4) {\begin{tabular}{@{}l@{}}
    \textcolor{impC}{\textgt{ReactionTrial}}（1回分）\\\hline
    \texttt{delayMs}：光るまでの待ち時間\\
    \texttt{reactionMs}：反応時間（\texttt{null} は失敗）
  \end{tabular}};
  \draw[arr] (se.west) -- node[lab,above,sloped]{\texttt{sessionId} で結ぶ} (ss.east);
  \draw[arr] (rr.west) -- node[lab,below,sloped]{\texttt{sessionId} で結ぶ} (ss.east);
  \draw[arr] (ss.south) -- node[lab,right]{\texttt{reward} として中に持つ} (ra.north);
  \draw[arr] (rr.south) -- node[lab,right]{\texttt{trials} に5回分} (rt.north);
\end{tikzpicture}
\end{center}
\figcap{図4　前後チェックの記録の形（\texttt{types/studyLog.ts}）。矢印は「どれとつながるか」}

\begin{codebox}{impA}{types/studyLog.ts}{1}
export const LOCAL_USER_ID = 'local-user-001';

export interface RewardAppraisal {
  interest: number; // 1-5 面白さ
  importance: number; // 1-5 重要度
  hasDeadline: boolean;
  deadlineUrgency?: number; // 1-5、hasDeadline=trueの時のみ
}

export interface StudySession {
  id: string;
  userId: string;
  subject?: string;
  startedAt: string; // ISO
  endedAt?: string; // ISO
  durationMinutes: number;
  reward: RewardAppraisal | null;
  subjectivePreId?: string;
  subjectivePostId?: string;
  reactionPreId?: string;
  reactionPostId?: string;
  /** 将来の「予定」欄用。現時点ではUIなし、欠損許容。 */
  plannedFor?: string; // YYYY-MM-DD
}

export type SubjectiveTiming = 'pre' | 'post';

export interface SubjectiveEntry {
  id: string;
  userId: string;
  sessionId?: string;
  timing: SubjectiveTiming;
  score: number; // 1-5（疲労感・高いほど疲れている）
  note?: string;
  timestamp: string; // ISO
}

export interface ReactionTrial {
  delayMs: number;
  reactionMs: number | null; // null = 反応なし/フライング
}
\end{codebox}

1行目の \texttt{LOCAL\_USER\_ID} は、全部の記録に入れている仮のユーザー番号。
ログインしても、記録はこの同じ番号で端末に保存される。
14行目の \texttt{// ISO} は、日時を \texttt{"2026-11-17T05:30:00.000Z"} のような文字列で持つという意味（末尾の \texttt{Z} は世界標準時）。

疲労の予定は、もう1つのファイルで形を決めている。

\begin{codebox}{impB}{types/fatigue.ts}{1}
export type FatigueAxis = 'cognitive' | 'physical';
export type FatigueWeight = 1 | 2 | 5;
export type FatigueLevel = 'ok' | 'caution' | 'danger';

export interface ScheduleTag {
  id: string;
  name: string;
  axis: FatigueAxis;
  weight: FatigueWeight;
  emoji: string;
}

/**
 * 正規化された予定エントリ。
 * v3: source='manual'（手動入力）のみ。
 * v4: source='calendar'（Googleカレンダー等）を追加予定。
 * calcFatigueScore() はどちらのソースも同一ロジックで処理できる。
 */
export interface ScheduleEntry {
  id: string;
  tagId: string;
  date: string; // YYYY-MM-DD
  durationMinutes: number;
  source: 'manual' | 'calendar';
  title?: string; // v4でカレンダーイベント名を上書き
  externalId?: string; // v4でカレンダーイベントID
}
\end{codebox}

\texttt{ScheduleTag} は「部活」「授業」のような予定の種類で、\texttt{axis} が身体（\texttt{physical}）か頭（\texttt{cognitive}）か、
\texttt{weight} がその重さ（1・2・5のどれか）。\texttt{ScheduleEntry} は「その種類の予定を、何日に何分やったか」。
21行目の \texttt{tagId} で \texttt{ScheduleTag} とつながる。

最後に、学習時間そのものの形。これは \texttt{TimeContext.tsx} の中にある。

\begin{codebox}{impA}{context/TimeContext.tsx}{6}
const STORAGE_KEY = 'timeLogData';

type SavedData = {
  totalMinutes: number;
  history: Record<string, number>;
};
\end{codebox}

\texttt{totalMinutes} は累計の分数、\texttt{history} は「日付 → その日の分数」の辞書。
たとえば \texttt{\{ "2026-11-16": 50, "2026-11-17": 45 \}} のようになる。

\textgt{端末の保存箱（AsyncStorage）の一覧}

\begin{center}
\begin{tblr}{
  colspec={Q[l,wd=3.6cm] X[l] Q[l,wd=3.6cm] Q[c,wd=1.6cm]},
  width=\linewidth,
  row{1}={bg=kkdark, fg=white, font=\small\gtfamily},
  hlines={black!15},
  cells={font=\small},
}
キー（箱の名前） & 中身 & 担当 & 重要度 \\
\SetRow{bg=impAbg} \texttt{timeLogData} & 累計の分数と、日ごとの分数 & \texttt{TimeContext.tsx} & \lvA \\
\SetRow{bg=impAbg} \texttt{studyLogData\_v1} & 学習1回ごとの記録・前後の疲労感・反応時間 & \texttt{StudyLogContext.tsx} & \lvA \\
\SetRow{bg=impBbg} \texttt{fatigueData\_v3} & 予定の種類と、日ごとの予定 & \texttt{FatigueContext.tsx} & \lvB \\
\SetRow{bg=impCbg} \texttt{appSettings} & テーマと換算比率 & \texttt{SettingsContext.tsx} & \lvC \\
\SetRow{bg=impCbg} \texttt{authGuestMode} & ゲストで使っているかどうか & \texttt{AuthContext.tsx} & \lvC \\
\end{tblr}
\end{center}

\subsection{1つの機能を縦に追う：タイマー → 保存 → ホーム表示}

ここがこのノートの中心。「タイマーで開始して、停止すると、ホームの『今日の学習時間』が増える」までを、
図5の番号①〜⑬の順に、元のコードで追う。

\begin{center}
\begin{tikzpicture}[font=\scriptsize,
  head/.style={draw=#1, fill=#1!10, line width=0.8pt, rounded corners=2pt,
    minimum width=2.6cm, minimum height=1.35cm, align=center, anchor=north, inner sep=2pt},
  msg/.style={-{Stealth[length=2mm]}, line width=0.7pt, draw=black!75},
  ret/.style={-{Stealth[length=2mm]}, line width=0.7pt, draw=black!60, dashed},
  ml/.style={fill=white, inner sep=1.2pt, align=center, above}]
  \node[head=impC] (hu) at (0,0) {\footnotesize\textgt{あなた}};
  \node[head=impA] (ht) at (2.95,0) {\footnotesize\textgt{タイマー画面}\\\texttt{timer.tsx}};
  \node[head=impB] (hm) at (5.9,0) {\footnotesize\textgt{前後チェック}\\\texttt{PreSessionModal}\\\texttt{PostSessionModal}};
  \node[head=impA] (hc) at (8.85,0) {\footnotesize\textgt{共有データ}\\\texttt{StudyLogContext}\\\texttt{TimeContext}};
  \node[head=impB] (hs) at (11.8,0) {\footnotesize\textgt{端末の保存箱}\\AsyncStorage};
  \node[head=impA] (hh) at (14.75,0) {\footnotesize\textgt{ホーム画面}\\\texttt{index.tsx}};
  \foreach \n in {hu,ht,hm,hc,hs,hh} { \draw[black!30, dashed] (\n.south) -- (\n.south |- 0,-12.6); }
  % 開始
  \draw[msg] (0,-1.9) -- node[ml]{①「開始」を押す} (2.95,-1.9);
  \draw[msg] (2.95,-2.75) -- node[ml]{② 前チェックを開く} (5.9,-2.75);
  \draw[msg] (5.9,-3.6) -- node[ml]{③ 入力を渡す（\texttt{onDone}）} (2.95,-3.6);
  \draw[msg] (2.95,-4.45) -- node[ml]{④ \texttt{startSession()} で記録を作る} (8.85,-4.45);
  \draw[msg] (8.85,-5.3) -- node[ml]{⑤ \texttt{studyLogData\_v1} に保存} (11.8,-5.3);
  \node[draw=impA, fill=impAbg, rounded corners=2pt, align=center, inner sep=2pt] at (2.95,-6.15)
    {⑥ 1秒ごとに \texttt{seconds} を +1};
  % 停止
  \draw[msg] (0,-7.0) -- node[ml]{⑦「停止」を押す} (2.95,-7.0);
  \draw[msg] (2.95,-7.85) -- node[ml]{⑧ 後チェックを開く} (5.9,-7.85);
  \draw[msg] (5.9,-8.7) -- node[ml]{⑨ 入力を渡す（\texttt{onDone}）} (2.95,-8.7);
  \draw[msg] (2.95,-9.55) -- node[ml]{⑩ \texttt{addMinutes(分)} と \texttt{finishSession()}} (8.85,-9.55);
  \draw[msg] (8.85,-10.4) -- node[ml]{⑪ \texttt{timeLogData} に保存} (11.8,-10.4);
  \draw[msg] (8.85,-11.25) -- node[ml]{⑫ 値が変わったので描き直す（\texttt{useTime()}）} (14.75,-11.25);
  \draw[ret] (14.75,-12.1) -- node[ml]{⑬「今日の学習時間」が増えて見える} (0,-12.1);
\end{tikzpicture}
\end{center}
\figcap{図5　タイマーを使ってから、ホームの数字が増えるまで（実線＝呼び出し、点線＝あなたに見える結果）}

\subsubsection{①② 「開始」を押すと、前チェックが開く}

\begin{codebox}{impA}{app/(tabs)/timer.tsx}{132}
        <View style={styles.buttonArea}>
          {status === 'idle' && (
            <TouchableOpacity
              style={[styles.btn, styles.btnPrimary]}
              onPress={() => setPreModalVisible(true)}
              activeOpacity={0.8}
            >
              <Text style={styles.btnTextLight}>開始</Text>
            </TouchableOpacity>
          )}
\end{codebox}

133行目の \texttt{status === 'idle' \&\&} は「状態が \texttt{idle}（止まっている）のときだけ、このボタンを表示する」という書き方。
136行目の \texttt{onPress} が押したときの処理で、\texttt{setPreModalVisible(true)} で前チェックの画面を「見える」にする。
\textgt{このボタンはまだタイマーを動かさない。}動かすのは③のあと。

タイマー画面が覚えている値（state）は、関数の最初にまとめて書いてある。

\begin{codebox}{impA}{app/(tabs)/timer.tsx}{23}
export default function TimerScreen() {
  const colorScheme = useColorScheme();
  const isDark = colorScheme === 'dark';

  const { addMinutes } = useTime();
  const { startSession, finishSession, addSubjectiveEntry, addReactionResult } = useStudyLogContext();
  const [status, setStatus] = useState<TimerStatus>('idle');
  const [seconds, setSeconds] = useState(0);
  const intervalRef = useRef<ReturnType<typeof setInterval> | null>(null);
  const [preModalVisible, setPreModalVisible] = useState(false);
  const [postModalVisible, setPostModalVisible] = useState(false);
  const currentSessionIdRef = useRef<string | null>(null);
\end{codebox}

27〜28行目で、共有データ（Context）から使う関数を取り出している。
29〜33行目が state で、\texttt{status} は \texttt{'idle'}（停止）・\texttt{'running'}（計測中）・\texttt{'paused'}（一時停止）のどれか。
34行目の \texttt{currentSessionIdRef} は「今計っている学習の番号」で、画面の見た目に関係しないので state ではなく \texttt{useRef} にしている。

\subsubsection{③④⑤ 前チェックの入力を受け取り、記録を作って保存する}

\begin{codebox}{impA}{app/(tabs)/timer.tsx}{61}
  const handlePreSessionDone: ComponentProps<typeof PreSessionModal>['onDone'] = input => {
    const sessionId = startSession({ subject: input.subject, reward: input.reward });
    if (input.subjectiveScore != null) {
      addSubjectiveEntry({
        sessionId,
        timing: 'pre',
        score: input.subjectiveScore,
        note: input.subjectiveNote,
      });
    }
    if (input.reaction) {
      addReactionResult({
        sessionId,
        timing: 'pre',
        trials: input.reaction.trials,
        avgReactionMs: input.reaction.avgReactionMs,
        accuracy: input.reaction.accuracy,
      });
    }
    currentSessionIdRef.current = sessionId;
    setPreModalVisible(false);
    setStatus('running');
  };
\end{codebox}

前チェックで「完了」を押すと、この関数が呼ばれる（③）。
62行目で \texttt{startSession} を呼び、新しい学習の記録を作ってその番号をもらう（④）。
63〜79行目は、疲労感や反応時間を入力した場合だけ、それも記録する（入力は任意なので \texttt{if} で確かめている）。
82行目の \texttt{setStatus('running')} で、ここで初めてタイマーが動き出す。

\begin{codebox}{impA}{context/StudyLogContext.tsx}{107}
  const startSession = useCallback(
    (input: { subject?: string; reward: RewardAppraisal | null; plannedFor?: string }) => {
      const id = genId();
      const session: StudySession = {
        id,
        userId: LOCAL_USER_ID,
        subject: input.subject,
        startedAt: new Date().toISOString(),
        durationMinutes: 0,
        reward: input.reward,
        plannedFor: input.plannedFor,
      };
      setSessions(prev => {
        const next = [...prev, session];
        persist({ sessions: next });
        return next;
      });
      return id;
    },
    [persist]
  );
\end{codebox}

110〜118行目で、図4の \texttt{StudySession} の形どおりに記録を作る。この時点では \texttt{durationMinutes} は 0。
120行目で、今までの記録の配列の後ろに新しい記録を足した配列を作り、
121行目の \texttt{persist} で AsyncStorage の \texttt{studyLogData\_v1} に保存する（⑤）。

\begin{pyaside}
  120行目の \texttt{[.{}.{}.prev, session]} は \texttt{[*prev, session]}、
  つまり「\texttt{prev} の中身を全部並べて、最後に \texttt{session} を足した新しいリスト」。
  \texttt{prev.append(session)} のように元のリストを書き換えないのは、
  React が「新しい値になった」と気づいて描き直すため。
\end{pyaside}

\subsubsection{⑥ 1秒ごとに数える}

\begin{codebox}{impA}{app/(tabs)/timer.tsx}{36}
  useEffect(() => {
    if (status === 'running') {
      intervalRef.current = setInterval(() => {
        setSeconds(prev => prev + 1);
      }, 1000);
    } else {
      if (intervalRef.current) {
        clearInterval(intervalRef.current);
        intervalRef.current = null;
      }
    }

    return () => {
      if (intervalRef.current) {
        clearInterval(intervalRef.current);
        intervalRef.current = null;
      }
    };
  }, [status]);
\end{codebox}

54行目の \texttt{[status]} により、この処理は \texttt{status} が変わったときだけ動く。
\texttt{'running'} になったら、38〜40行目の \texttt{setInterval} で 1000 ミリ秒（1秒）ごとに \texttt{seconds} を 1 増やす。
一時停止や停止で \texttt{'running'} でなくなったら、43行目の \texttt{clearInterval} で止める。
48〜53行目の \texttt{return} の中身は「後片付け」で、画面を離れたときにもタイマーを止める。

\begin{pyaside}
  開発ログ V1 によると、\texttt{setInterval}／\texttt{clearInterval} は \texttt{threading.Timer} にあたる。\evLog
  39行目の \texttt{prev => prev + 1} は \texttt{lambda prev: prev + 1}。
\end{pyaside}

\subsubsection{⑦⑧⑨⑩ 「停止」を押し、後チェックのあとで分数を足す}

\begin{codebox}{impA}{app/(tabs)/timer.tsx}{56}
  const handleStop = () => {
    setStatus('idle');
    setPostModalVisible(true);
  };
\end{codebox}

停止ボタンは、タイマーを止めて（57行目）後チェックを開く（58行目）だけ。まだ何も保存しない。

\begin{codebox}{impA}{app/(tabs)/timer.tsx}{85}
  const handlePostSessionDone: ComponentProps<typeof PostSessionModal>['onDone'] = input => {
    const elapsedMinutes = Math.round(seconds / 60);
    addMinutes(elapsedMinutes);

    const sessionId = currentSessionIdRef.current;
    if (sessionId) {
      finishSession(sessionId, elapsedMinutes);
\end{codebox}

後チェックで「完了」を押すと、この関数が呼ばれる（⑨）。
86行目で秒を分に直す。\texttt{Math.round} は四捨五入なので、29秒までは 0 分、30秒からは 1 分になる。
87行目の \texttt{addMinutes} で \texttt{TimeContext} に分数を足し、
91行目の \texttt{finishSession} で \texttt{StudySession} に終了時刻と分数を書き込む（⑩）。
このあと92〜108行目は、③のときと同じように \texttt{'post'}（後）の疲労感と反応時間を記録する。

\begin{codebox}{impA}{app/(tabs)/timer.tsx}{109}
    }
    currentSessionIdRef.current = null;
    setPostModalVisible(false);
    setSeconds(0);
  };
\end{codebox}

最後に、学習の番号を空にして、後チェックを閉じ、秒を 0 に戻す。

\subsubsection{⑪ TimeContext が分数を足して保存する}

\begin{codebox}{impA}{context/TimeContext.tsx}{59}
  useEffect(() => {
    if (!loaded) return;
    const data: SavedData = { totalMinutes, history };
    AsyncStorage.setItem(STORAGE_KEY, JSON.stringify(data)).catch(() => {});
  }, [totalMinutes, history, loaded]);

  const addMinutes = (minutes: number) => {
    if (minutes <= 0) return;
    const today = todayKey();
    setTotalMinutes(prev => prev + minutes);
    setHistory(prev => ({
      ...prev,
      [today]: (prev[today] ?? 0) + minutes,
    }));
  };
\end{codebox}

\textgt{読む順番は下から。}⑩で呼ばれるのは 65行目の \texttt{addMinutes}。
\begin{itemize}
  \item 66行目：0 分以下なら何もしない。
  \item 67行目：今日の日付 \texttt{"2026-11-17"} のような文字列を作る（\texttt{utils/date.ts} の関数）。
  \item 68行目：累計に足す。
  \item 69〜72行目：日ごとの辞書の、今日の分に足す。今日の分がまだなければ 0 から足す（\texttt{?? 0}）。
\end{itemize}
68〜72行目で state が変わると、59〜63行目の \texttt{useEffect} が「\texttt{totalMinutes} と \texttt{history} が変わった」と気づいて動き、
JSON の文字列にして \texttt{timeLogData} に保存する（⑪）。
\textgt{「足す」処理と「保存する」処理が分かれていて、保存は state の変化をきっかけに自動で起きる}のがポイント。

\begin{pyaside}
  62行目の \texttt{JSON.stringify(data)} は \texttt{json.dumps(data)}。
  70〜71行目の \texttt{\{ .{}.{}.prev, [today]: … \}} は \texttt{\{**prev, today: …\}}。
  71行目の \texttt{prev[today] ?? 0} は \texttt{prev.get(today, 0)} とほぼ同じ。
\end{pyaside}

\subsubsection{⑫⑬ ホーム画面が描き直される}

\begin{codebox}{impA}{context/TimeContext.tsx}{29}
export function TimeProvider({ children }: { children: ReactNode }) {
  const [totalMinutes, setTotalMinutes] = useState(0);
  const [history, setHistory] = useState<Record<string, number>>({});
  const [loaded, setLoaded] = useState(false);

  const todayMinutes = history[todayKey()] ?? 0;
\end{codebox}

\texttt{todayMinutes}（今日の分数）は保存されている値ではなく、34行目で \texttt{history} から毎回計算している。
\texttt{history} が変わると \texttt{TimeProvider} が描き直され、この値も新しくなる。
そして、\texttt{useTime()} でこの値を使っているすべての画面に「値が変わった」と伝わる（⑫）。

\begin{codebox}{impA}{app/(tabs)/index.tsx}{61}
export default function HomeScreen() {
  const colorScheme = useColorScheme();
  const isDark = colorScheme === 'dark';
  const { totalMinutes, todayMinutes } = useTime();
  const { fatigueState } = useFatigue();
\end{codebox}

\begin{codebox}{impA}{app/(tabs)/index.tsx}{170}
          <View style={styles.statItem}>
            <ThemedText style={styles.statLabel}>今日の学習時間</ThemedText>
            <View style={styles.statValueRow}>
              <ThemedText style={styles.statValueBig}>{todayMinutes}</ThemedText>
              <ThemedText style={styles.statUnit}>分</ThemedText>
            </View>
          </View>
\end{codebox}

ホーム画面は 64行目で \texttt{useTime()} から値を受け取り、173行目の \texttt{\{todayMinutes\}} の場所に表示する。
ホーム画面のコードには「記録を読み込む」処理も「描き直す」処理も書いていない。
\textgt{Context の値が変わると、それを使っている画面は React が自動で描き直す}から（⑬）。
キャラクターの成長（累計 60 分・300 分・600 分・1500 分で次の姿になる）も、同じ \texttt{totalMinutes} から決まる。

\begin{kkbadgepoint}[要点]{この流れから分かる、アプリの基本の形}
  画面（\texttt{timer.tsx}・\texttt{index.tsx}）は、記録を自分では持たない。
  記録を変えたいときは Context の関数（\texttt{addMinutes} など）を呼び、表示したいときは Context の値（\texttt{todayMinutes} など）を受け取るだけ。
  保存は Context の中の \texttt{useEffect} が、値が変わるたびに自動で行う。
  だから「タイマー画面で足す」と「ホーム画面に出る」が、互いのコードを知らなくてもつながる。
  ほかの機能（疲労・設定）も、この同じ形でできている。
\end{kkbadgepoint}

\subsection{細部：出てきたライブラリと関数}

この章のコードに出てきたものだけをまとめる。「どこから来たか」で色を分けた。

\begin{longtblr}[entry=none,label=none]{
  colspec={Q[l,wd=3.3cm] Q[l,wd=2.4cm] X[l] Q[l,wd=3.4cm]},
  width=\linewidth,
  rowhead=1,
  row{1}={bg=kkdark, fg=white, font=\small\gtfamily},
  hlines={black!15},
  cells={font=\small},
}
名前 & どこから来たか & 何をするか & Pythonでいうと \\
\SetRow{bg=impAbg} \texttt{useState} & React & 覚えておく値と、それを変える関数を作る & インスタンス変数＋setter \\
\SetRow{bg=impAbg} \texttt{useEffect} & React & 指定した値が変わったときに処理を動かす & 値を見張るコールバック \\
\SetRow{bg=impAbg} \texttt{createContext}・\texttt{useContext} & React & 共有データの箱を作る・箱から取り出す & Flask の \texttt{g} \\
\SetRow{bg=impBbg} \texttt{useRef} & React & 描き直しても消えない入れ物 & インスタンス変数 \\
\SetRow{bg=impBbg} \texttt{useCallback} & React & 関数を毎回作り直さずに覚えておく & なし（速さのための工夫） \\
\SetRow{bg=impAbg} \texttt{AsyncStorage.getItem}・\texttt{setItem} & ライブラリ（async-storage） & 端末の保存箱から読む・書く & ファイルの読み書き \\
\SetRow{bg=impBbg} \texttt{JSON.stringify}・\texttt{JSON.parse} & JavaScript 標準 & データと文字列を相互に変換 & \texttt{json.dumps}・\texttt{json.loads} \\
\SetRow{bg=impBbg} \texttt{setInterval}・\texttt{clearInterval} & JavaScript 標準 & 一定時間ごとに処理を動かす・止める & \texttt{threading.Timer} \\
\SetRow{bg=impBbg} \texttt{Math.round} & JavaScript 標準 & 四捨五入 & \texttt{round}（ただし 2.5 は Python だと 2、JS だと 3） \\
\SetRow{bg=impBbg} \texttt{new Date().toISOString()} & JavaScript 標準 & 今の日時を世界標準時の文字列にする & \texttt{datetime.now(timezone.utc).isoformat()} \\
\SetRow{bg=impAbg} \texttt{todayKey}・\texttt{localDateKey} & 自作（\texttt{utils/date.ts}） & 端末の時刻で \texttt{"YYYY-MM-DD"} を作る & \texttt{date.today().isoformat()} \\
\SetRow{bg=impBbg} \texttt{genId} & 自作（\texttt{StudyLogContext.tsx}） & 時刻と乱数から記録の番号を作る & \texttt{uuid4()} の簡易版 \\
\SetRow{bg=impCbg} \texttt{View}・\texttt{Text}・\texttt{TouchableOpacity} & React Native & 箱・文字・押せる部品 & tkinter の \texttt{Frame}・\texttt{Label}・\texttt{Button} \\
\SetRow{bg=impCbg} \texttt{StyleSheet.create} & React Native & 見た目の設定をまとめる & tkinter の \texttt{configure()}（開発ログ V1） \\
\SetRow{bg=impCbg} \texttt{SafeAreaView} & ライブラリ（safe-area-context） & 画面の切り欠きを避けて表示する & なし \\
\SetRow{bg=impCbg} \texttt{=>}（アロー関数） & JavaScript の書き方 & 短く書ける関数 & \texttt{lambda} \\
\SetRow{bg=impCbg} \texttt{.{}.{}.}（スプレッド構文） & JavaScript の書き方 & 配列や辞書の中身を広げる & \texttt{*list}・\texttt{**dict} \\
\SetRow{bg=impCbg} \texttt{??} & JavaScript の書き方 & 左が空（\texttt{null}／\texttt{undefined}）なら右を使う & \texttt{x if x is not None else 0} \\
\SetRow{bg=impCbg} \texttt{条件 \&\& 部品} & JavaScript の書き方 & 条件が正しいときだけ部品を表示 & \texttt{if 条件: 表示} \\
\end{longtblr}

